\documentclass{rebuttal}
\renewcommand{\locallang}{PL}
\input{preamble.tex}

\addbibresource{rebuttal.bib}

\title{Odpowiedź na pytania Recenzenta}
\author{Kacper Kania}
\date{}


\begin{document}
\maketitle

\responsedelimiter

\paragraph{Pytanie.}
  Mimo krótkiej wzmianki na temat roli i udziału współautorów w publikacjach,
  warto byłoby jakby Doktorant w punktach wymienił swój indywidualny wkład w
  każdą publikację.
\paragraph{Odpowiedź.}
  Zgadzam się z Recenzentem, że wartość pracy doktorskiej zostałaby podniesiona,
  gdyby zawierała dodatkową informację o wkładzie autorów w publikacje, tak jak
  jest przyjęte w przypadku publikacji w czasopismach wydawanych przez IEEE. W
  każdej z pierwszoautorskich publikacji miałem znaczący wpływ na rozwój
  projektu. To uwzględnia zadania takie jak:
  \begin{itemize}
    \item określenie problemu badawczego oraz pytań z tym związanych,
    \item postawienie hipotezy badawczej,
    \item wykonywanie eksperymentów testujących zadaną hipotezę,
    \item pisanie kodu metody oraz eksperymentów,
    \item budowanie aparatu matematycznego formalizującego działanie metody,
    \item pisanie publikacji oraz udział w poprawianiu oraz nadzorowaniu aplikowania
          poprawek do publikacji.
  \end{itemize}
  Ostatni punkt wiąże się też z budowaniem
  materiałów promocyjnych metod, np. \url{https://conerf.github.io},
  \url{https://blendfields.github.io}, \url{https://lumigauss.github.io}. W
  przypadku pracy dotyczącej LumiGaussa, moim wkładem było nadzorowanie
  projektu, konsultacje pomysłów, wykonanie jednego z eksperymentów, oraz
  redakcja Rysunków w finalnej publikacji oraz formalizmu matematycznego
  zawartego w publikacji.

  Pragnę przy tym zaznaczyć wpływ, jaki mieli współautorzy na mój rozwój oraz
  postęp projektów---od \textit{brainstormingu} na temat idei projektu oraz jego
  kierunku, przez konsultacje związane z opracowaniem metody i budowaniem sieci
  kontaktów pomocnych przy rozwoju projeku, po pomoc w redagowaniu publikacji
  oraz nadzorowaniu rozwoju moich umiejętności związanych z umiejętnościami w
  prowadzeniu badań naukowych i ich opisywaniem.

  \responsedelimiter

\paragraph{Pytanie.}
  Zaproponowana sieć neuronowa dla systemu CoNeRF to stosunkowo prosta w swojej
  strukturze sieć typu MLP, co jest w pewnym sensie powieleniem pomysłu z
  oryginalnych publikacji dotyczących metody NeRF. Ciekawe jednak czy nie lepsza
  byłaby tu sieć grafowa? Tego typu podejście umożliwiłoby uwzględnienie
  przestrzennej struktury pikseli już w samej architekturze sieci. Czy Doktorant
  rozważał użycie architektur sieci innych niż MLP? To samo pytanie można też
  powtórzyć w przypadku metody BlendFields.

\paragraph{Odpowiedź.} Rzeczywiście, kodowanie sąsiedztwa czy to przez sieć
  grafową~\cite{ververas2024sags}, konkatenację współrzędnych sąsiadujących
  punktów przy użyciu $k$-NN~\cite{dupont2022from}, czy użycie
  Transformera~\cite{joshi2025transformers,zhuang2023diffusion} przyczynia się
  do zwiększenia jakości rekonstrukcji w ujęciu generatywnym. Okazuje się
  jednak, że zastosowanie wystarczająco ekspresywnej sieci MLP jest
  wystarczające w zadaniu rekonstrukcji. Oba stwierdzenia potwierdziły
  wcześniejsze moje badania wykonywane przy każdym z prezentowanych projektów.

  Pragnę zaznaczyć również, że operacje w sieciach MLP są wysoce zoptymalizowane
  i ,,tanie'' w ujęciu wykorzystania pamięci. To pozwalało nam efektywnie
  wykonywać eksperymenty i iterować w projekcie po różnych pomysłach. Analiza
  wpływu bardziej skomplikowanych modeli jest ciekawym kierunkiem, który
  potencjalnie pozwoliłby osiągnięcie jeszcze lepszych wyników osiąganych przez
  zaproponowane metody.

  Celem tego projektu było również zaproponowanie podejścia, które pozwala na
  osiągnięcie kontroli predykcji przy rekonstrukcji 3D, niezależnie od samej
  architektury wykonującej tę predukcję. Aby zminimializować ilość stopni
  swobody wpływających na jakość rekonstrukcji, zastosowałem architekturę
  sprawdzoną we wcześniejszych
  badaniach~\cite{mildenhall2020nerf,park2021nerfies}.

  W przypadku BlendFields, zrealizowałem rozszerzenie metody (nazwane
  InstantBlendFields) z bardziej zaawansowaną architekturą, które pozwalało na
  szybsze uczenie modelu. Materiały promocyjne podejścia można obejrzeć pod tym
  linkiem: \url{https://blendfields.github.io/instantblendfields/}.

  \responsedelimiter

\paragraph{Pytanie.}
  Rozdział 2.3.3, Wzór (19) - Czy użyty tu symbol mnożenia to mnożenie w sensie
  Khatri-Rao?

\paragraph{Odpowiedź.}
  W tej i pozostałych pracach nadużywamy typowej notacji. Takie nadużycie w
  wielu przypadkach pozwala na uniknięcie wizualnego zanieczyszczenia formalizmu
  matematycznego, które utrudniałoby zrozumienie pracy. Dodatkowo, zostało
  nieformalnie przyjęte, że operator w postaci symbolu $\odot$ oznacza wykonanie
  mnożenia \textit{element-wise} (element po elemencie, lub elementów o tych
  samych indeksach w obu macierzach). W środowiskach typu Matlab taka operacja
  jest zazwyczaj zapisywana jako~\texttt{*.}, gdzie ten zapis nie jest
  powszechnie stosowany w publikacjach.

  \responsedelimiter

\paragraph{Pytanie.}
  Rozdział 2.3.4 - skoro w sieci kodujemy również pozycje punktów (,,windowed
  positional encoding'') to nie lepiej byłoby rozważyć architekturę transformera
  wizyjnego, który już w naturalny dla siebie sposób dokonuje tego typu
  kodowania?
\paragraph{Odpowiedź.}
  Architektura sieci wykorzystująca same sieci MLP były w wielu przypadkach
  wystarczające oraz tańsze pod względem wykorzystania pamięci na GPU dostępnych
  w tamtym momencie. Dodatkowo, prace analizujące korzyści płynące z
  wykorzystania warunkowanych transformerów zamiast sieci MLP do rekonstrukcji
  3D~\cite{rebain2023attention} pojawi się równolegle do naszej pracy.

  Zgadzam się, że wykorzystanie transformera w tym przypadku mogłoby zwiększyć
  jakość rekonstrukcji. Jednak samo kodowanie pozycyjne w ich sytuacji nie
  podlega tak bezpośredniej kontroli jak to ma w przypadku \textit{windowed
    positional encoding}~\cite{park2021nerfies}.

  \responsedelimiter

\paragraph{Pytanie.}
  Rozdział 2.4.1 -- co to znaczy ,,we set the attributes to be synchronized for
  the training split, and desynchronized for the test split''? Jak dokonuje się
  tego typu ,,synchronizacji'' atrybutów danych?

\paragraph{Odpowiedź.}
  Celem wykorzystania syntetycznego zbioru było pokazanie, że nasza metoda może
  działać w skrajnych przypadkach, kiedy kontrolowalne atrybuty obiektów są
  zsynchronizowane, natomiast w czasie testowanie podlegają losowym, ciągłym
  zmianą. Na przykładzie użytego syntetycznego zbioru danych, synchronizacja
  atrybutów oznacza jednoczesną zmianę kolorów obiektów w dokładnie w ten sam
  sposób. Natomiast w czasie testu, model przedstawiający czajnik zmienia kolor
  na zielony, natomiast model królika na czerwony (czyli zmiany są
  zdesynchronizowane). Przykładowa wizualizacja tej synchronizacji jest
  przedstawiona
  \href{https://conerf.github.io/videos/synthetic-data/train.mp4}{tutaj},
  natomiast desynchronizacji pod
  \href{https://conerf.github.io/videos/synthetic-data/valid.mp4}{tym linkiem}.
  Pragnę zaznaczyć, że ten zbiór został wygenerowany sztucznie za pomocą
  programu Blender.

  \responsedelimiter

\paragraph{Pytanie.}
  Przy prezentowaniu metod warto byłoby pokazać przypadki, dla których one
  jednak nie działają lub działają poniżej oczekiwań. Czy jest wiele takich
  przypadków? Kiedy i dla jakich danych?
\paragraph{Odpowiedź.}
  Tę odpowiedź pragnę podzielić na cztery części, po jednej dla każdej z metod.

  \textbf{CoNeRF} wykazywał szczególne niesatysfakcjonujące w trzech
  przypadkach. Pierwszym z nich są drobne szczegóły obiektów, które chcemy
  modelować ze względu na ograniczoną ilość danych---aby drobne, ruchome obiekty
  były poprawnie rekonstruowane, dane trenujące musiałyby być bardzo wysokiej
  jakości, w jednolitym oświetleniu z wielu ujęć dokoła. Dodatkowo,
  niewystarczające oznaczenia danych mogą powodować skorelowanie wpływu wartości
  atrybutów na wyjście, co objawia się modyfikacją wielu obszarów jednocześnie
  na wyjściu mimo zmiany tylko pojedynczego atrybutu. Wszystkie te przypadki
  zostały przedstawione na Rysunku 16.

  \textbf{Blendfields} dziedziczy problemy podejść opartych na dokładnej
  kalibracji kamer. W trakcie eksperymentów zauważyliśmy, że niedokładne
  spasowanie modelu twarzy powodowało efekt \textit{ghostingu}, gdzie można było
  zauważyć wiele instancji osoby nakładających się niedokładnie ze sobą.
  Dodatkowo, niski kontrast zmarszczek wynikających ze zniekształceń objawiał
  się niską wartości funkcji straty, co powodowało niedokładne dopasowanie
  modelu do danych.

  \textbf{LumiGauss} ze względu na wykorzystanie harmonik sferycznych niskiego
  stopnia nie był w stanie modelować ostrych cieni wynikających z kierunkowego
  oświetlenia, np. słońca w bezchmurny dzień. Jednocześnie wykorzystanie
  wyższych stopni powodowało błędne dopasowanie do danych i migotanie przy
  ciągłej rotacji mapy oświetlenia.

  \textbf{CLoG} natomiast miał dwa znaczące ograniczenia wynikające z jego
  zdefiniowania. Po pierwsze, nie był w stanie modelować scen otwartych, w
  których scena nie ma określonych granic. Po drugie, jakość rekonstrukcji była
  niższa od standardowego podejścia Gaussian Splatting w scenach o dużej
  złożoności (np. scena z trawnikiem, gdzie pojedyncze źdźbła są widoczne na
  zdjęciach). To ograniczenie wynikało z użycia reprezentacji w postaci
  dwuwymiarowej mapy, gdzie mapowanie otwartej reprezentacji lub wielu
  szczegółów jest ograniczone.

  \responsedelimiter

\paragraph{Pytanie.}
  Rozdział 2.4.4 -- ,,our controllability losses do not interfere with the
  rendering quality'' -- jak zostało to zmierzone?
\paragraph{Odpowiedź.}
  Głównym celem zastosowanych składowych funkcji w CoNeRFie straty było
  nauczenie sieci dwóch rzeczy---poprawnej rekonstrukcji obrazów po
  zaaplikowaniu równania renderującego oraz rekonstrukcji adnotacji masek, które
  były dostępne dla niektórych ze zdjęć. By dodatkowo wymusić gładkość
  rekonstrukcji, czyli aby dla dwóch oznaczeń czasowych $t_1, t_2, t_2 - t_1 \to
    0$ ich rekonstrukcje były maksymalnie podobne\footnote{W pracy powołuję się na
    warunek Lipschitza, czyli $\|f(t_2) - f(t_1)\| \leq \|t_2 - t_1\|$, gdzie
    $f(t)$ jest funkcją rekonstruującą obraz dla oznaczenia czasowego $t$.}. Celem
  tej pracy były jedynie dodanie możliwości kontroli wyjścia bez wpływu na
  jakość tej predykcji. W tym celu wykonałem ablację, gdzie mierzyłem jak
  dodanie kolejnych elementów funkcji straty wpływa na zadanie rekonstrukcji. W
  Tabeli~4 przedstawiłem wyniki i pokazałem, że przy zadaniu interpolacji
  widoków wyniki nie zmieniają się w porównaniu do zwykłego NeRFa, który jest
  warunkowany wyuczalnym, wielowymiarowym wektorem~\cite{mildenhall2020nerf}. W
  tym kontekście przedstawiłem wniosek, że zdefiniowane funkcje straty
  odpowiadający za dodanie kontroli predykcji nie wpływają na jakość
  renderowania.

  \responsedelimiter

\paragraph{Pytanie.}
  W przypadku prezentowanych metod warto byłoby żeby Doktorant umieścił linki do
  przykładowych sekwencji ukazujących ich właściwości, np. w postaci wideo.
  Kilka znalazłem, ale przez przypadek.
\paragraph{Odpowiedź.}
  Zgadzam się z Recenzentem, że linki do stron z materiałami powinny być
  umieszczone w pracy. Pragnę je przytoczyć tutaj:
  \begin{itemize}
    \item CoNeRF: \url{https://conerf.github.io/}
    \item BlendFields: \url{https://blendfields.github.io/}
    \item LumiGauss: \url{https://lumigauss.github.io/}
  \end{itemize}
  Materiały dotyczący CLoG nie zostały upublicznione.
  Można je znaleźć spakowane w pliku \texttt{zip} pod
  \href{https://drive.google.com/file/d/1ViaYDDo-x5DjfMBDzO1bjK4UDX1qEomv/view?
    usp=sharing}{tym linkiem}.

  \responsedelimiter

\paragraph{Pytanie.}
  Rozdział 3.2.4, Rozdział 4.4.2 i wiele innych miejsc -- dziwne użycie słowa
  ,,concurrent''. Być może chodziło o ,,competitive''?
\paragraph{Odpowiedź.}
  Publikacje znajdujące się w pracy doktorskiej były umieszczane z zawartością
  taką, jaką posiadały w momencie publikacji lub wysyłania do recenzji. To
  oznacza też, że w wielu przypadkach odnosiłem się do prac równoległych, czyli
  takich, które były publikowane w podobnym momencie, w którym pisałem daną
  publikację. Ze względu na chęć zachowania oryginalności treści publikacji, nie
  redagowałem ich przy umieszczaniu ich w pracy doktorskiej.

  \responsedelimiter

\paragraph{Pytanie.}
  Rozdział 5.3.5 -- nie do końca jasne jest uzasadnienie naukowe takiego, a nie
  innego wyboru architektury sieci MLP.
\paragraph{Odpowiedź.}
  Zastosowana architektura była dobrana empirycznie w czasie wczesnych
  eksperymentów. Została ona oparta na jednej z pierwszych prac dotyczącej
  NeRFów~\cite{mildenhall2020nerf}. Jej dobór został podyktowany chęcią łatwego
  dopasowania sieci do danych treningowych przy jednoczesnym zachowanie
  ograniczenia pamięci wykorzystywanej wtedy karty graficznej A100 z 40GB VRAM.

  \responsedelimiter

\paragraph{Pytanie.}
  Obecnie obserwujemy bardzo dynamiczny rozwój badań oraz metod dotyczących
  połączenia, czy też sterowania procesami generatywnymi z a pomocą wielkich
  modeli językowych LLM. Co prawda w końcowym rozdziale 6.2 Autor coś wspomina o
  rozwiązaniach typu ,,text-prompted'', jednakże bardzo brakuje mi powiązania
  tego nurtu z tym co Doktorant prezentuje w swojej rozprawie.
\paragraph{Odpowiedź.}
  Zgadzam się z recenzentem, że wykorzystanie modeli językowych do budowy
  generatywnych podejść rekonstrukcji 3D jest bardzo ciekawym kierunkiem. W tej
  rozprawie skupiłem się na podejściu rekonstrukcji \textit{per} scena ze
  względu na ograniczenia obliczeniowe. Dodatkowo znaczący rozwój tych modeli
  pojawił się w późniejszym okresie, kiedy prace na prezentowanymi publikacjami
  były już w zaawansowanym stadium. Pragnę zaznaczyć jak bardzo rozwinęła się ta
  klasa modeli i przytoczyć kilka przykładowych kategorii oraz prac, które
  zostały opracowane w tym czasie:
  \begin{itemize}
    \item generowanie pojedynczych obiektów \cite{tang2024dreamgaussian},
    \item generowanie całych scen \cite{bai2025360},
    \item segmentacja obiektów w przestrzeni 3D \cite{kobayashi2022decomposing},
    \item edytowanie pojedynczych obiektów \cite{haque2023instruct}.
  \end{itemize}

  \responsedelimiter

  \printbibliography[title={Bibliografia}]

\end{document}